% answers_short.tex —— 1984 年·全国卷（文科） · 简短答案速查
% 由 tools/make_short_answers.py 自 answers.tex 的 \gksolhead/\gkanswer 提取，
% 并与 1984/1984_数学真题_文_答案要点.md 逐题比对后生成，勿手改（改源 MD 后重跑）。

\begin{gknotice}
\textbf{使用说明}：本节只给最终答案便于快速核对，完整推导见解析版。答案为编写者独立推导，\textbf{非官方标准答案}。
\end{gknotice}

\gkpart{题目}
\begin{itemize}[leftmargin=2.6em]
  \item[\textbf{1.}] C
  \item[\textbf{2.}] D
  \item[\textbf{3.}] A
  \item[\textbf{4.}] C
  \item[\textbf{5.}] A
  \item[\textbf{6.}] $\dfrac32<x<2$
  \item[\textbf{7.}] $\dfrac4\pi$ 或 $\dfrac8\pi$
  \item[\textbf{8.}] $m=0$，$x=-\dfrac12$
  \item[\textbf{9.}] $1$
  \item[\textbf{10.}] $6!\cdot\mathrm A_7^4$
  \item[\textbf{11.}] 抛物线 $y^2=-4x$——顶点在原点，开口向左，对称轴为 $x$ 轴，焦点为 $(-1,0)$，准线为直线 $x=1$
  \item[\textbf{12.}] 函数 $y=\dfrac{1}{(x+1)^2}$ 的图象由两支组成——全在 $x$ 轴上方，以直线 $x=-1$ 和 $y=0$ 为渐近线，关于直线 $x=-1$ 对称；在 $(-\infty,-1)$ 上单调递增、在 $(-1,+\infty)$ 上单调递减（越靠近 $x=-1$ 越向上无限伸展），过点 $(0,1)$ 与 $(-2,1)$
  \item[\textbf{13.}] 见解析版
  \item[\textbf{14.}] $\sin(\alpha+\beta)\sin(\alpha-\beta)$
  \item[\textbf{15.}] $\dfrac{64\sqrt3}{3}\ \mathrm{cm}^3$
  \item[\textbf{16.}] 从 1993 年开始
  \item[\textbf{17.}] (1) $C_1$ 的中心为 $(0,0)$、焦点为 $(\pm2\sqrt{10},0)$；$C_2$ 的中心为 $(3,0)$、焦点为 $(3\pm4\sqrt2,0)$；(2) $(x-14)^2+y^2=125$ 或 $(x+1)^2+y^2=20$
\end{itemize}
